% Transcription — fonds Alexandre Grothendieck, shelfmark 5, batch 1 (pages 1-20).
% Established from the facsimile archives/batches/5/batch-01.pdf (a mirror of
% https://grothendieck.umontpellier.fr/5.pdf), per the skill
% transcribe-grothendieck. The batch file carries no cover sheet: PDF page N is
% archivists' page N. Checked against the pencilled numbers bottom-left ("2",
% "4", "5", "6", "7", "8", "17", "18" read on the corresponding leaves).
%
% Pass: Fable 5.1 (claude-fable-5-1), 2026-09-19 — first pass, unchecked against the pages by a human.
%
% What the batch is. Three short runs in his hand, wrapped in covers, on the
% versos of typescripts that are not his; the folder's inventory title says
% "cristaux et cohomologie de De Rham : généralités", and that is what the
% hand carries. Page 1 is a blank tan wrapper (only the show-through of page
% 2) and gets no \page. Page 2, pencil, is a single leaf: the diagram of the
% topoi attached to X — Zariski, "conn", P-conn, stratifying, crystalline,
% P-cris, infinitesimal — with the equivalences noted on the arrows, then
% three remarks on what the arrows to and from X_zar give and on the Chern
% class of an invertible Module in H^2 of the crystalline topoi. Pages 4 to 8,
% ink, paginated 4-8 by the archivists and not by him, are one run on the
% universal artinian W-algebra with nu-nilpotent divided powers: the ring
% W^{(nu)}, the auxiliary A^{(nu)}, the normal form of its elements (two
% lemmas, a corollary), the quotients W^{(nu,r)}, and on page 8 the
% independent computation with the Lambda_i and the exponent bound
% sum a_i p^i >= p^{nu+1}. Page 10 is a half-leaf (landscape) in ink on the
% same subject. Page 12, pencil, is a single leaf on de Rham cohomology and
% the crystalline site. Page 14 is a tan cover in his hand, "Frobenius", for
% the run that follows: pages 15 to 19, pencil, "Général nonsens sur
% Frobenius" — the topos Q(alpha, beta) of pairs (E, F_E) with F_E : alpha^*E
% -> beta^*E for two morphisms alpha, beta : X' => X, its universal property,
% its functoriality, and the question whether omega^* commutes with the R^i
% u_*; page 19 carries a single line closing page 18. Page 20 is a second tan
% cover in his hand naming the run that begins in batch 2; it gets no \page
% and its words are not made a heading here, because the run it names has no
% leaf in this batch.
%
% Skipped, without description, as the skill requires: page 3 (typescript, not
% his), page 9 (typescript, not his), page 11 (typescript, not his), page 13
% (a list of page-and-line corrections to a typescript, in his hand, but not
% mathematics in hand).
%
% Notation fixed once here. The topoi of page 2 are set X with a roman
% subscript (X_{zar}, X_{conn}, X_{cris}, X_{inf}, X_{strat}) and the bold P of
% the leaf is \mathbb{P} (X_{\mathbb{P}-conn}, X_{\mathbb{P}-cris}). On pages
% 4-8 his capital lambda and capital theta are \Lambda_i and \Theta_i, his
% small lambda is \lambda_i; the dotted Theta of the corollary on page 6 and
% of the lemma on page 7 (the image in the quotient) is \dot{\Theta}_i; the
% binomial coefficient he writes C^p_i is kept as such; the divided power
% gamma^{p^r} is \gamma^{p^{r}}; the underlined structure sheaf and the
% underlined L of page 2 are \underline{\mathcal{O}} and \underline{L}. On pages
% 15-19 the topos Q(alpha, beta) is \mathcal{Q}(\alpha,\beta) (his Q is a
% looped capital), the forgetful morphism is \omega, the Frobenius-type
% morphisms of a topos are F_X, F_Y, the homomorphisms attached to a sheaf are
% F_E, F_G, and the base-change maps are c_E and c^G as he indexes them.
%
% Readings flagged and not settled: on page 2 the underlined phrase after
% "isomorphismes de localisation", read "zariskienne"; the class c'(L), which
% may be c^1(L); the label on the right vertical arrow of the diagram, read
% "équiv. si X/p lisse". On page 4 the parenthesis after "algèbres sur W",
% read "(Cohen)"; the script letter before Lambda in "s'envoie dans rΛ", read
% as a radical. On page 5 the crossed symbol before w_a in a), and the middle
% of the binomial expansion, where a diagonal stroke cancels part of the
% right-hand side. On page 16 the exact words after "Pour prouver que Q est
% un topos"; on page 17 the words between "d'ailleurs les foncteurs" and the
% setting X' = X.
\documentclass[11pt,a4paper]{article}
% Le préambule commun à toutes les transcriptions du fonds Grothendieck.
%
% Il tient en une idée : **l'incertitude se compose, elle ne se résout pas.**
% Une transcription qui lisse un mot douteux a détruit la seule chose pour
% laquelle on la faisait — un lecteur doit pouvoir savoir, à chaque endroit,
% ce qui était sur le feuillet et ce qui a été ajouté. Les six macros qui
% suivent sont donc l'appareil critique, et non de la décoration.
%
% Les mêmes commandes sont reconnues par `scripts/render.mjs`, qui produit la
% vue de lecture du site. Une macro ajoutée ici doit l'être aussi là-bas,
% faute de quoi le rendu échouera bruyamment — ce qui est voulu : un rendu
% silencieusement faux est pire qu'un rendu absent.

\ProvidesPackage{grothendieck}[2026/08/08 Transcriptions du fonds Grothendieck]

\usepackage{amsmath,amssymb,amsthm}
\usepackage{mathrsfs}
\usepackage{stmaryrd}
% Les diagrammes commutatifs sont partout dans ce fonds : c'est la langue
% même de la géométrie algébrique de Grothendieck, et une transcription qui
% les rendrait en prose ne transcrirait rien.
\usepackage{tikz-cd}
% Français par défaut — c'est la langue des feuillets — mais l'anglais est
% chargé aussi : la traduction et le résumé sont des documents anglais, et la
% typographie française (espaces avant les deux-points) y serait une faute.
% Ces documents basculent par \selectlanguage{english} en tête de corps.
\usepackage[english,french]{babel}
\usepackage{fontspec}
\usepackage{geometry}
\geometry{a4paper,margin=25mm,marginparwidth=28mm}
\usepackage{marginnote}
\usepackage{xcolor}
% ulem, not soul. `soul`'s \st and \ul are built on a character-by-character
% scan that XeLaTeX's font machinery breaks: the document fails with an
% undefined control sequence rather than degrading. ulem does the same two jobs
% and is Unicode-engine safe. `normalem` stops it hijacking \emph, which the
% transcriptions use for Grothendieck's own emphasis.
\usepackage[normalem]{ulem}
\usepackage{hyperref}
\hypersetup{colorlinks=true,linkcolor=black,urlcolor=black,pdfborder={0 0 0}}

% --- Métadonnées du lot -----------------------------------------------------
% Reprises telles quelles par le script de rendu, qui les affiche en tête de
% la vue de lecture. Elles fixent ce que le fichier prétend couvrir : sans
% elles, une transcription flotte sans point d'attache dans un fonds de seize
% mille feuillets.

\newcommand{\folder}[1]{\gdef\grothFolder{#1}}
\newcommand{\batch}[1]{\gdef\grothBatch{#1}}
\newcommand{\leaves}[2]{\gdef\grothFirst{#1}\gdef\grothLast{#2}}
\newcommand{\foldertitle}[1]{\gdef\grothFoldertitle{#1}}
\gdef\grothFolder{?}\gdef\grothBatch{?}\gdef\grothFirst{?}\gdef\grothLast{?}\gdef\grothFoldertitle{}

% --- Le résumé --------------------------------------------------------------
%
% Il ouvre la lecture modernisée, sous le titre « Résumé », et il est composé
% en petit. Ce n'est pas un ornement : c'est le seul passage du document qui
% ne soit pas des mathématiques, et un lecteur doit pouvoir le reconnaître
% avant de l'avoir lu — puis le sauter s'il connaît déjà le sujet, ou s'y
% arrêter s'il ne le connaît pas. Le filet qui le suit marque la reprise du
% corps.

\newenvironment{resume}
  {\small\setlength{\parskip}{0.45em}}
  {\par\medskip\hrule\medskip}

% --- Mots-clés ---------------------------------------------------------------
%
% En anglais, en clôture du résumé de la lecture modernisée : le vocabulaire
% moderne sous lequel chercher ce que les feuillets construisent. Le manifeste
% du site (`npm run manifest`) les extrait pour étiqueter la cote entière —
% cette ligne est la source unique des tags, il n'y a pas de fichier à côté.
\newcommand{\keywords}[1]{\par\smallskip\noindent
  {\footnotesize\textsc{Keywords} --- \textit{#1}}\par}

% --- L'appareil critique ----------------------------------------------------

% Le numéro de feuillet, en marge. C'est l'ancre qui permet de régler un
% désaccord sur une formule en regardant *un* feuillet plutôt qu'une chemise
% de six cents. Le site s'en sert aussi pour tourner les pages du fac-similé
% quand on fait défiler la transcription.
\newcommand{\leaf}[1]{%
  \marginnote{\footnotesize\color{gray!70}\textsf{#1}}%
  \label{leaf:#1}\ignorespaces}

% Illisible : on ne devine pas. Un mot inventé qui se lit comme les autres est
% le pire résultat possible.
\newcommand{\ill}{\textcolor{red!60!black}{[\,\ldots\,]}}

% Lecture douteuse : le mot est donné, et signalé comme incertain.
\newcommand{\uncertain}[1]{\uline{#1}}

% Ajout de l'éditeur — une lettre manquante, un mot sous-entendu.
\newcommand{\add}[1]{\textcolor{blue!50!black}{[#1]}}

% Biffé par Grothendieck lui-même. Ce qu'il a rayé fait partie du document :
% une rature dit souvent où la pensée a changé de direction.
\newcommand{\struck}[1]{\sout{#1}}

% Note du transcripteur, sur l'état matériel du feuillet.
\newcommand{\note}[1]{\par\smallskip{\small\color{gray!60!black}\textit{[#1]}}\par\smallskip}

% Note marginale de Grothendieck, distincte du corps du texte.
\newcommand{\marginal}[1]{\par\smallskip\noindent\hspace{1em}%
  {\small\color{gray!60!black}#1}\par\smallskip}

% --- Titre ------------------------------------------------------------------

\AtBeginDocument{%
  \begin{center}
    {\large\bfseries Fonds Alexandre Grothendieck --- cote n\textsuperscript{o}~\grothFolder}\\[2pt]
    {\itshape\grothFoldertitle}\\[4pt]
    {\small feuillets \grothFirst--\grothLast\ (lot \grothBatch)}\\[2pt]
    {\footnotesize\color{gray!60!black}Transcription établie d'après le fac-similé numérisé
     par l'Université de Montpellier.\\
     Les lectures incertaines sont soulignées, les illisibles notées [\,\ldots\,],
     les ajouts entre crochets.}
  \end{center}
  \vspace{1em}}

\endinput


\folder{5}
\batch{1}
\pages{1}{20}
\dating{1965-1969}
\watermark{Édition de démonstration}
\foldertitle{Cristaux et cohomologie de De Rham : généralités, correspondances Deligne et Berthelot (1965-1969) : notes manuscrites (s.d.), lettre (1965).}

\begin{document}

\section*{Les topos attachés à $X$}

\page{2}

\begin{tikzcd}[column sep=small, row sep=small, nodes={font=\scriptsize}]
 & & X_{\mathrm{zar}} \arrow[dll] \arrow[drr] & & \\
X_{\mathrm{conn}} \arrow[rrrr, dashed] \arrow[drr, "\text{équiv. en car. } 0"] \arrow[dd, "\text{équiv. si } X/S \text{ lisse}"'] & & & & X_{\mathbb{P}\mathrm{-conn}} \arrow[dll, dashed] \arrow[dd, "\text{équiv. si } X/p \text{ lisse}"] \\
 & & X_{\mathrm{strat}} \arrow[dd] & & \\
X_{\mathrm{cris}} \arrow[rrrr, dashed] \arrow[drr] \arrow[ddrr] & & & & X_{\mathbb{P}\mathrm{-cris}} \arrow[dll, dashed] \arrow[ddll] \\
 & & X_{\mathrm{inf}} \arrow[d] & & \\
 & & X_{\mathrm{zar}} & &
\end{tikzcd}

\note{le diagramme occupe la moitié supérieure du feuillet ; les flèches en
tirets sont en tirets sur la page. Une accolade à droite du diagramme porte :
« morphismes de topos \emph{annelés} [$\dashrightarrow$ définis si $X$ de torsion
sur $\mathbb{Z}$] ». Le label de la flèche verticale de droite est lu « équiv.
si $X/p$ lisse », celui de gauche « équiv. si $X/S$ lisse ». Un trait fin, sans
pointe visible, descend en outre de $X_{\mathrm{zar}}$ (en haut) vers
$X_{\mathrm{inf}}$ ; il n'est pas rendu comme flèche}

\marginal{NB \quad $S$ sous-entendu dans notations.}

\marginal{détermination des \emph{points} de ces topos ?}

Les flèches à valeurs dans $X_{\mathrm{zar}}$ donnent lieu à des
\emph{isomorphismes de localisation \uncertain{zariskienne}}, resp. des suites
spectrales intéressantes.

Les flèches en provenance de $X_{\mathrm{zar}}$ donnent lieu à une suite exacte
\[
\cdot \longrightarrow J \longrightarrow \underline{\mathcal{O}}_{\mathcal{X}}
\longrightarrow \underline{\mathcal{O}}_{\mathcal{X}_0} \longrightarrow \cdot
\]
\marginal{dire que $\mathrm{Mod}(X_{\mathrm{zar}}) \xrightarrow{\;\approx\;}
\mathrm{Mod}(\underline{\mathcal{O}}_{\mathcal{X}_0})$}\note{la note est reliée
par un trait à la suite exacte}

Pour un Module inversible $\underline{L}$ sur $X_{\mathrm{zar}}$, on définit une
classe de Chern $c'(\underline{L})$\note{lu $c'$ ; peut-être $c^{1}$} dans
$H^2(X_{\mathrm{cris}})$ et $H^2(X_{\mathbb{P}\mathrm{-cris}})$, qui sont
compatibles si $X$ de torsion. Sur $X_{\mathrm{cris}}$, son annulation est NS
pour que $\underline{L}$ provienne d'un Module inversible sur
$X_{\mathrm{cris}}$ [donc dans le cas relatif lisse, \struck{d'un} qu'il admette
une connexion intégrable].

\section*{L'algèbre artinienne universelle à puissances divisées $\nu$-nilpotentes}

\page{4}

$W$ vecteurs de Witt sur corps parfait $k$.

$\nu \geq 1$ entier.

On s'intéresse à la catégorie des \struck{algèbres} anneaux artiniens
\struck{sur} \struck{augmentés vers} de corps rés. $k$, à idéal maximal muni de
puiss. divisées $\nu$-nilpotentes. On va construire un objet initial de cette
catégorie.

Pour ceci, on note que les $\Lambda$ \uncertain{sont} des algèbres sur $W$
(\uncertain{Cohen}), et que l'élément $\lambda_0 = p \in W$ s'envoie dans
$\mathfrak{r}\Lambda$\note{la lettre devant $\Lambda$ est lue comme un
$\mathfrak{r}$, le radical}. Donc $\Lambda$ reçoit
\[
W^{(\nu)} = W[\Lambda_0, \dots, \Lambda_\nu] \Big/ \Big(\Lambda_0^{p} - p!\,\Lambda_1,\;
\Lambda_1^{p} - \frac{p^2!}{(p!)^{p}}\,\Lambda_2,\; \dots,\;
\Lambda_{\nu-1}^{p} - \frac{p^{\nu}!}{(p^{\nu-1}!)^{p}}\,\Lambda_\nu,\;
\Lambda_\nu^{p} \;;\; \Lambda_0 - \struck{p}\lambda_0\Big)
\]
[\struck{$W$-algèbre \ill{}} avec idéal $J = (\Lambda_0, \dots, \Lambda_\nu)$ à
puiss. div. $\nu$-nilpotentes, \emph{universel} pour les algèbres sur $W$, avec
idéal à puiss divisées $\nu$-nilpotentes contenant l'image de $\lambda_0$]. NB.
$\Lambda_\nu = \gamma^{p^{\nu}}(\Lambda_0)$ ; on utilise le fait que dans $W$
les nbs premiers $\ell \neq p$ sont inversibles.

Or \struck{\ill{}} on voit tout de suite que $W^{(\nu)}$ est artinien
\add{l'idéal maximal étant précisément $J$}\note{ajout interlinéaire de sa
main} de corps résiduel $k$, \struck{(ceci est vrai avant d'ajouter la
relation $\Lambda_0 - \lambda_0$), \ill{}}

On va calculer explicitement $W^{(\nu)}$. Pour ceci, \struck{\ill{}} on commence
par l'anneau plus simple
\[
A^{(\nu)} = W[\Lambda_0, \dots, \Lambda_\nu] \big/ \big(\Lambda_0^{p} - p!\,\Lambda_1,\; \dots,\; \Lambda_\nu^{p}\big)
\]
(on oublie la relation $\Lambda_0 - \lambda_0$). On pose aussi
\[
\lambda_r = \frac{p^{p^{r}}}{p^{r}!} = \gamma^{p^{r}}_{\text{habituel}}(p)
\]
et on pose
\[
\Lambda_i = \lambda_i + \Theta_i \qquad 0 \leq i \leq \nu .
\]
Par abus de \ill{}, on désigne \ill{} par $\Lambda_i$, $\Theta_i$ les images de
$\Lambda_i$, $\Theta_i$ dans $A^{(\nu)}$. Alors $A^{(\nu)}$ a une

\page{5}

base sur $W$ formée des
\[
\Theta_0^{a_0} \cdots \Theta_\nu^{a_\nu}, \qquad 0 \leq a_i \leq p-1 \quad \text{pour } i = 0, \dots, \nu .
\]

Il faut maintenant diviser par l'idéal $I$ engendré par $\Theta_0$. On a le

\textbf{Lemme} \quad Pour que
$\displaystyle\sum_{\substack{a = (a_0, \dots, a_\nu) \\ 0 \leq a_i \leq p-1}} w_a\, \Theta_0^{a_0} \cdots \Theta_\nu^{a_\nu}$
soit dans l'idéal engendré par $\Theta_0$, il faut et il suffit que

a) pour tout $a = (a_0, \dots, a_\nu)$, si \add{$a \neq 0$ et si}\note{ajout
interlinéaire de sa main} $a_i$ est la première composante non nulle de $a$, on
ait \struck{\ill{}} $w_a \in p^{i} W$

b) le coefficient $w_0$ de $1$ soit dans $p^{c(\nu)} W$, où
\[
c(\nu) = \nu + p\,\big[\,p^{\nu} - (p^{\nu-1} + p^{\nu-2} + \dots + 1)\,\big]
\]

\textbf{Dém.} \quad \emph{Condition suffisante} \quad Résulte des relations
\[
\begin{cases}
p^{i}\,\Theta_i \in I \\
p^{c(\nu)} \in I
\end{cases}
\]
La première relation par récurrence sur $i \geq 0$. Si $i = \nu$\note{sic ; on
attend $i = 0$} \ill{} trivial, i.e. $\Theta_0 \in I$. Supposons \ill{}
$p^{i}\Theta_i \in I$ ($i \leq \nu - 1$) on utilise
\[
(\Theta_i + \lambda_i)^{p} = \frac{p^{i+1}!}{(p^{i}!)^{p}}\,(\Theta_{i+1} + \lambda_{i+1}) ,
\]
\ill{}
\[
\struck{\ill{}} \sum_{\alpha=1}^{p-1} C_p^{\alpha}\,\Theta_i^{\alpha}\lambda_i^{p-\alpha}
= \Big( -\lambda_i^{p} + \frac{p^{i+1}!}{(p^{i}!)^{p}}\,\struck{\Theta_{i+1}}
+ \frac{p^{i+1}!}{(p^{i}!)^{p}}\,\Theta_{i+1} \Big)
\]
\note{un trait diagonal barre une partie du membre de droite ; le terme barré
est lu $\Theta_{i+1}$ et devait être $\lambda_{i+1}$}
multipliant par $p^{i}$, \ill{} $p^{i}\Theta_i \in I$ donc $p^{i}\Theta_i^{\alpha} = 0$
si $\alpha \geq 1$ :
\[
p^{i}\,\frac{p^{i+1}!}{(p^{i}!)^{p}}\,\Theta_{i+1} = 0
\]
or on constate que $v_p\Big(\dfrac{p^{i+1}!}{(p^{i}!)^{p}}\Big) = 1$

\page{6}

donc les relations précédentes équivalent à
\[
p^{i+1}\,\Theta_{i+1} = 0 .
\]
On a donc prouvé les relations $p^{i}\Theta_i = 0$. Pour la dernière relation, on
écrit
\[
(\lambda_\nu + \Theta_\nu)^{p} = 0 \quad \text{i.e.} \quad
\lambda_\nu^{p} + \sum_{i=1}^{p} C_i^{p}\,\Theta_\nu^{i}\lambda_\nu^{p-i} = 0
\]
Multiplions par $p^{\nu}$ et utilisons que $p^{\nu}\Theta_\nu \in I$ donc
$p^{\nu}\Theta_\nu^{\alpha} \in I$ si $\alpha \geq 1$, on trouve
\[
p^{\nu}\lambda_\nu^{p} \in I
\]
Or \ill{}
\[
v_p(p^{\nu}\lambda_\nu^{p}) = \nu + p\,v_p(\lambda_\nu)
= \nu + p\Big(p^{\nu} - \frac{p^{\nu}-1}{p-1}\Big) = c(\nu)
\]
\note{sous $\lambda_\nu$ il rappelle $\lambda_\nu = p^{p^{\nu}}/p^{\nu}!$}

Pour prouver que la condition est nécessaire, il suffit de noter que l'ensemble
$I'$ formé des éléments satisfaisant les cond. a) b) données (qui satisfont
$I' \subset I$, $\Theta_0 \in I'$) est bien un \emph{idéal}, i.e. \struck{\ill{}}
(comme c'est un sous-$W$-module) est stable par multiplication par
$\Theta_0, \dots, \Theta_\nu$. La vérification est facile et laissée au lecteur.
\marginal{!}

\textbf{Corollaire} \quad Soit, pour tout $i$, $E_i$ un syst. de représentants
de $W/p^{i}W$. Alors un élément de $W^{(\nu)}$ s'écrit de façon unique sous la
forme
\[
\sum w_{a_1, \dots, a_\nu}\, \dot{\Theta}_1^{a_1} \cdots \dot{\Theta}_\nu^{a_\nu}
\qquad 0 \leq a_i \leq p-1 \quad (1 \leq i \leq \nu)
\]
avec
\[
\begin{cases}
w_{a_1, \dots, a_\nu} \in E_i \quad \text{si $a_i$ est la première coordonnée non nulle de } (a_1, \dots, a_\nu) \\
w_0 \in E_{c(\nu)}
\end{cases}
\]

\page{7}

\note{le feuillet reprend la seconde ligne de l'accolade, $w_0 \in E_{c(\nu)}$,
qui clôt le corollaire de la page 6}

\textbf{Lemme} \quad Soit $1 \leq r \leq \nu$, et soit $W^{(\nu, r)}$ le quotient
de $W^{(\nu)}$ par $\dot{\Theta}_1, \dots, \dot{\Theta}_r$. \struck{Alors les}
Soit \struck{$A$} $I_{\nu, r}$ l'idéal de $A^{(\nu)}$ engendré par
$\Theta_0, \dots, \Theta_r$, de sorte que $A^{(\nu)}/I_{\nu, r} \simeq W^{(\nu, r)}$.
Pour que $\sum w_a\,\Theta_0^{a_0} \cdots \Theta_\nu^{a_\nu}$ soit dans
$I_{\nu, r}$, il f. et s. que

a) pour tout $a = (a_0, \dots, a_\nu) \neq 0$, \struck{\ill{}}
\struck{$(\alpha,\ 1 \leq \alpha \leq \nu - r,\ \ill{})$} si $a_i$ est la
première composante non nulle et si $i > r$, on ait \struck{\ill{}}
$w_a \in p^{i-r} W$

b) on ait $w_0 \in p^{c(\nu) - r} W$

\textbf{Cor.} \quad Forme normale d'un élément de $W^{(\nu, r)}$ :
\[
\sum w_{a_{r+1}, \dots, a_\nu}\, \Theta_{r+1}^{a_{r+1}} \cdots \Theta_\nu^{a_\nu}
\]
avec
\[
\begin{cases}
w_{a_{r+1}, \dots, a_\nu} \in E_i\struck{\ill{}} \quad \text{si $a_{r+i}$ est le premier des $a_\alpha$ qui est } \neq 0 \\
w_0 \in E_{c(\nu) - r}
\end{cases}
\]

\page{8}

$W$ vecteurs de Witt

$\Lambda_0 \cdots \Lambda_\nu$ indéterminées, liées par
\[
\begin{cases}
\Lambda_i^{p} = \dfrac{p^{i+1}!}{(p^{i}!)^{p}}\,\Lambda_{i+1} & \text{si } 0 \leq i \leq \nu - 1 \\[6pt]
\Lambda_\nu^{p} = 0
\end{cases}
\]
Ceci implique
\[
\Lambda_0^{a_0} \cdots \Lambda_\nu^{a_\nu} = 0 \quad \text{si} \quad \sum a_i p^{i} \geq p^{\nu+1} \qquad (a_i \geq 0)
\]

\struck{Comme} \textbf{Dém} \struck{Récurrence sur les} $\sum a_i$, \ill{} si
$\sum a_i = 0$ \ill{}\note{deux lignes biffées et surchargées ; seule leur
armature se lit} \marginal{Si tous les $a_i$ sont $\leq p-1$, c'est ok, car
$\sum a_i p^{i} \leq p^{\nu+1}$ \ill{}. Supposons que $a_i \geq p$, donc
\ill{}}\note{la note est encadrée à droite, avec un renvoi vers le
$\sum a_i p^{i} \geq p^{\nu+1}$ de l'énoncé}

\textbf{Lemme}
\[
\Lambda_0^{a_0}\Lambda_1^{a_1} \cdots \Lambda_\nu^{a_\nu}
= c\,\Lambda_0^{b_0} \cdots \Lambda_{\nu-1}^{b_{\nu-1}}\Lambda_\nu^{b_\nu}, \qquad c \in W,
\]
avec $0 \leq b_i \leq p-1$ pour $0 \leq i \leq \nu - 1$, $b_\nu$ \uncertain{ce
qu'il faut}, et
\[
\sum_0^{\nu} a_i p^{i} = \sum_0^{\nu} b_i p^{i} .
\]
Pour le voir, on procède par récurrence \struck{\ill{} transformer
l'expression en une autre \ill{} tous les $a_\alpha \leq p-1$ \ill{}
$\alpha \leq i-1$ ($i \leq \nu - 1$) \ill{} $\sum a_\alpha p^{\alpha}$ \ill{}}
\struck{Donc} \ill{}\note{trois lignes biffées}. Donc si \ill{} $a_i \geq p$,
on écrit $a_i = p q_i + a'_i$, d'où
\[
\Lambda_i^{a_i} = (\Lambda_i^{p})^{q_i}\,\Lambda_i^{a'_i}
= (c_i\Lambda_{i+1})^{q_i}\,\Lambda_i^{a'_i}
= c\,\Lambda_{i+1}^{q_i}\,\Lambda_i^{a'_i}
\]
\note{les exposants de cette ligne sont écrits petit et se lisent mal ; ils
sont restitués d'après $a_i = pq_i + a'_i$}
donc
\[
\Lambda_0^{a_0} \cdots \Lambda_\nu^{a_\nu}
= c\,\Lambda_0^{a_0} \cdots \Lambda_i^{a'_i}\,\Lambda_{i+1}^{a_{i+1} + q_i}\,\Lambda_{i+2}^{a_{i+2}} \cdots ,
\]
et évidemment on a
\[
\sum a_i p^{i} = \sum a'_i p^{i} ,
\]
car \struck{$a'_i p^{i} +$}
\[
a'_i p^{i} + (a_{i+1} + q_i)\,p^{i+1} = \underbrace{(a'_i + q_i p)}_{a_i}\,p^{i} + a_{i+1}\,p^{i+1}
\]
OK.\note{les deux dernières lignes sont soulignées d'un double trait} Or si
$\sum b_i p^{i} \geq p^{\nu+1}$, alors \ill{} cela donne $b_\nu \geq p$, donc
$\Lambda_\nu^{b_\nu}$ est un multiple de $\Lambda_\nu^{p}$, or
$\Lambda_\nu^{p} = 0$. OK.

\page{10}

\note{demi-feuillet écrit en travers ; il calcule $v_p(n!)$ en base $p$. Il
écrit $v_p(p)$, $v_p(p^{2})$, … pour $v_p(p!)$, $v_p(p^{2}!)$, …
dans la colonne de gauche ; la transcription garde sa graphie}
\[
v_p(n!) = a_\nu\,v_p(p^{\nu}!) + \dots + a_{\mu+1}\,v_p(p^{\mu+1}!)
\]
\[
- v_p(p^{\mu+1}!) + (\mu+1) + (p-1)\big[\,v_p(p^{\mu}!) + \dots + v_p(p)\,\big]
\]
\note{le second membre est souligné d'une accolade portant un $p^{\mu}$ biffé ;
le lien entre les deux lignes n'est pas écrit}

\[
v_p(1) = 0
\]
\[
v_p(p) = 1 + (p-1)\cdot 0 = 1
\]
\[
v_p(p^{2}) = 2 + (p-1)\cdot 1 = p+1
\]
\[
v_p(p^{3}) = 3 + \struck{(p-1)}(p+2)\struck{+1} = p^{2} + p + 1
\]
\[
v_p(p^{4}) = 4 + (p-1)\big[p^{2} + 2p + 3\big] = 4 + p^{3} + 2p^{2} + 3p - p^{2} - 2p - 3
= \struck{1+}\,p^{3} + p^{2} + p + 1
\]
\[
v_p(p^{\nu}) = p^{\nu-1} + p^{\nu-2} + \dots + 1 = \frac{p^{\nu} - 1}{p - 1}
\]
\[
v_p(p^{\nu+1}) = \nu + 1 + (p-1)\big[p^{\nu-1} + 2p^{\nu-2} + 3p^{\nu-3} + \dots + \nu p^{0}\big]
\]
\[
= \nu + 1 + p^{\nu} + 2p^{\nu-1} + 3p^{\nu-2} + \dots + \nu p
- p^{\nu-1} - 2p^{\nu-2} - \dots - (\nu-1)p - \nu
\]
\[
= p^{\nu} + p^{\nu-1} + \dots + p + 1
\]
\[
\sum a_\nu\,(p^{\nu-1} + p^{\nu-2} + \dots + 1)
\]
\[
a_1 + a_2(p+1) + a_3(p^{2} + p + 1) + a_4(p^{3} + p^{2} + p + 1) + \dots
\]
\note{cette ligne est encadrée ; sous elle, $a_1 + p a_2 \to + p^{2} a_3$ et
un $p$ isolé}
\[
v_p(n!)(p-1) = a_1(p-1) + a_2(p^{2}-1) + \dots + a_\nu(p^{\nu} - 1)
= n - a_0 - a_1 - \dots - a_\nu
\]
\[
\boxed{\;v_p(n!) = \frac{n - \sum a_i}{p - 1}\;}
\]
\note{la formule est encadrée deux fois}

\page{12}

\[
v_p\Big(\frac{\pi^{n}}{n!}\Big) \geq 0 \quad \text{pour tout } n \ \ill{}
\]
\[
v_p\Big(\frac{\pi^{n}}{n!}\Big) \to +\infty \quad \text{pour } n \to +\infty
\]
\struck{$(p-1)$}
\[
v_p\Big(\frac{\pi^{n}}{n!}\Big) = \struck{(p-1)}\; n\,v_p(\pi) - \frac{n - \mathrm{chif}(n)}{p-1} \geq 0
\]
\[
\Updownarrow \qquad v_p(\pi) \geq \frac{1 - \frac{\mathrm{chif}(n)}{n}}{p-1}
\]
\struck{$v_p(\pi)$} \quad $\mathrm{chif}(n)/n \to 0$ \quad
\struck{$\frac{\mathrm{chif}(n)}{n}$ fonction décroissante} \quad
$\mathrm{chif}(n) \leq 1$\note{sic ; sans doute $\mathrm{chif}(n)/n \leq 1$}

de $n$, tous $\leq 1$\note{fin d'une ligne dont le début est perdu sous un
trait}
\[
n = a_0 + a_1 p + \dots + a_\nu p^{\nu}
\]
\[
\mathrm{chif}(n) = a_0 + a_1 + \dots + a_\nu
\]
si $a_\nu \neq 0$, on a \ill{}
\[
\mathrm{chif}(n) \leq (p-1)\nu\struck{\ill{}}
\]
\[
n \geq p^{\nu}
\]
\[
\frac{\mathrm{chif}(n)}{n} \leq \frac{(p-1)\nu}{p^{\nu}} \to 0
\]

donc
\[
v_p\Big(\frac{\pi^{n}}{n!}\Big) \geq 0 \ \text{pour tout } n
\iff v_p\Big(\frac{\pi^{n}}{n!}\Big) > 0 \ \text{pour tout } n
\iff v_p(\pi) \geq \frac{1}{p-1} \quad \text{i.e.} \quad \pi^{p-1} \in pV
\]
\note{les deux membres extrêmes sont encadrés}
\[
(p-1)\,v_p\Big(\frac{\pi^{n}}{n!}\Big) = n\big[(p-1)v_p(\pi) - 1\big] + \mathrm{chif}(n)
\]
si $v_p(\pi) = \frac{1}{p-1}$, c'est $\mathrm{chif}(n)$, qui ne tend pas vers
$+\infty$ ; si $v_p(\pi) > \frac{1}{p-1}$, on trouve $\underline{\alpha n +
\mathrm{chif}(n)}$ \add{$\alpha > 0$}, qui tend vers $\infty$
\[
n\Big(\alpha + \frac{\mathrm{chif}(n)}{n}\Big) \sim \alpha n
\]
Donc \struck{condition}
\[
v_p\Big(\frac{\pi^{n}}{n!}\Big) \to +\infty \iff v_p(\pi) > \frac{1}{p-1}
\quad \text{i.e.} \quad p = \text{unit}\cdot\pi^{r}, \ \text{avec } r < p-1
\]
\note{les deux membres sont encadrés}

\section*{Frobenius}

\note{le titre est de sa main, sur la chemise qui précède, page 14 ; la
chemise elle-même ne porte rien d'autre}

\page{15}

\textbf{Général nonsens sur Frobenius}

1) Soit

\begin{tikzcd}
X \arrow[r, "F_X"] \arrow[d, "g"'] & X \arrow[d, "g"] \\
Y \arrow[r, "F_Y"] & Y
\end{tikzcd}

un diagramme commutatif, \struck{\ill{}} de topos. Si $E$ est un faisceau
d'ens sur $X$, on a \struck{(sous des conditions \ill{})} \struck{\ill{}} un
homomorphisme de changement de base
\[
(1.1) \qquad c_E : F_Y^{*}(g_*(E)) \longrightarrow g_*(F_X^{*}(E))
\]
Si donc on se donne un homomorphisme
\[
(1.2) \qquad F_E : F_X^{*}(E) \longrightarrow E
\]
on en conclut
\[
g_*(F_E) : g_* F_X^{*} E \longrightarrow g_*(E)
\]
d'où un homomorphisme composé
\[
(1.3) \qquad \struck{c_E}\; g_*(F_E) \circ c_E : F_Y^{*}(g_*(E)) \longrightarrow g_*(E)
\]
\struck{En d'autres termes, on \ill{}} Si $E$ faisceau de \ill{} et (1.2) hom.
de tels, alors (1.3) aussi, car (1.1) l'est. Faits analogues pour $F$ dans une
catégorie \ill{}.

De même, considérons un faisceau $G$ sur $Y$, d'où un isom
\[
(1.4) \qquad c^{G} : F_X^{*}(g^{*}(G)) \xrightarrow{\;\sim\;} g^{*}(F_Y^{*}(G))
\]
d'où, si on se donne un homomorphisme
\[
(1.5) \qquad F_G : F_Y^{*}(G) \longrightarrow G
\]
d'où
\[
g^{*}(F_G) : g^{*}F_Y^{*}(G) \longrightarrow g^{*}(G) ,
\]
un homomorphisme composé
\[
(1.6) \qquad g^{*}(F_G) \circ c^{G} : F_X^{*}(g^{*}(G)) \longrightarrow g^{*}(G)
\]

2) Soient
\[
\alpha, \beta : X' \rightrightarrows X
\]
deux morphismes de topos. \struck{Alors} Considérons la catégorie
$\mathcal{Q}(\alpha, \beta)$ dont les objets sont les couples $(E, F_E)$,
\struck{$E \in \mathrm{Ob}\,X$,}
\[
E \in \mathrm{Ob}\,X, \qquad F_E : \alpha^{*}(E) \longrightarrow \beta^{*}(E)
\]
[resp. $F_E$ un \emph{isom}]. \struck{\ill{} $X'$ \ill{}} On a un foncteur
« oubli » $\omega : \mathcal{Q}(\alpha, \beta) \longrightarrow X$,
$\omega(E, F_E) = E$.

\textbf{Prop.} \quad $\mathcal{Q}(\alpha, \beta)$ est un topos, et $\omega$
est le foncteur image inverse d'un morphisme de topos

\page{16}

\[
\omega : X \longrightarrow \mathcal{Q}(\alpha, \beta)
\]

Considérons les composés $\omega\alpha$, $\omega\beta$. Alors on a un
homomorphisme canonique (resp. isom. can.]
\[
F : \omega\beta \longrightarrow \omega\alpha
\]
\add{dans} $\underline{\mathrm{Homtop}}(X, \mathcal{Q}(\alpha, \beta))$\note{sic ;
on attend $X'$ à la place de $X$}.

Pour tout topos $Y$, considérons la catégorie
$\underline{\mathrm{Homtop}}(X, \alpha, \beta ; Y)$ des \struck{foncteurs}
couples $(\varphi, f)$ d'un morphisme de topos $\varphi : X \to Y$ et d'un
\[
f : \varphi\beta \longrightarrow \varphi\alpha
\]
[qui se \ill{} à l'aide de $\underline{\mathrm{Hom}}\struck{\mathrm{top}}(X, Y)
\rightrightarrows \underline{\mathrm{Homtop}}(X', Y)$]\note{le crochet est
serré dans la marge droite ; le verbe entre « qui se » et « à l'aide de » ne
se lit pas}. \uncertain{Alors} le foncteur évident
\[
g \longmapsto (g \circ \omega,\; g * F) \;:\;
\underline{\mathrm{Homtop}}(\mathcal{Q}(\alpha, \beta), Y) \longrightarrow
\underline{\mathrm{Homtop}}(X, \alpha, \beta ; Y)
\]

\page{17}

\textbf{Ce foncteur est une équivalence de catégories.}

Pour prouver que $\mathcal{Q}$ est un topos, le \ill{} est \ill{} on voit qu'il
admet une petite famille génératrice (il est évident que $\mathcal{Q}$ satisfait
les conditions de Giraud a) b) c), et que \ill{} exact \ill{} et commute aux
\emph{lim}). Dans le cas resp., $\mathcal{Q}$ s'interprète comme la
catégorie des \emph{sections cartésiennes} d'un topos \emph{fibré}
$\mathbb{F}$ \ill{} sur un topos \ill{}\note{il dessine ici deux points $x'$,
$x$ avec deux flèches $\alpha$, $\beta$ de $x'$ vers $x$},
$\mathbb{F}_{x'} = X'$, $\mathbb{F}_x = X$, \add{associés aux} changements de
base $\alpha^{*}$, $\beta^{*}$, et \ill{} appliquer SGA 4 I 9.25. Dans le cas
non resp., on considère \struck{topos} \ill{} catégorie fibrée $\mathbb{F}$
\ill{} sur \ill{}\note{un diagramme dans la marge : $x' \to x$ par $u$, $x'
\to x'_1$ par $f$, $x'_1 \to x$ par $v$, avec la mention « associés aux
changements de base $u$ », le tout encadré ; « à vérifier » dans la marge de
gauche} que les arguments de SGA 4 I 9.21 et suivants donnent le résultat
…

Prenons le cas où $X' = X$, $\beta = \mathrm{id}$, donc où
$\mathcal{Q}(\alpha, \beta) = \struck{\ill{}}$ la catégorie des couples
$(E, F_E)$, $F_E : \alpha^{*}(E) \to E$.\note{le paragraphe est encadré}

La \struck{catégorie} topos $\mathcal{Q}(\alpha, \beta)$ dépend
« fonctoriellement » de $\alpha, \beta : X' \rightrightarrows X$, de façon
précise, si on a un diagramme

\begin{tikzcd}
X' \arrow[r, "\alpha"] \arrow[d, "u_1"'] & X \arrow[d, "u_0"] \\
Y' \arrow[r, "\lambda"] & Y
\end{tikzcd}

\note{les deux flèches horizontales sont doubles sur la page, $\alpha, \beta$
en haut et $\lambda, \mu$ en bas ; le rendu n'en trace qu'une par ligne}

\page{18}

commutatif à isom. près, on a \ill{} (p. ex. : l'un des \ill{} morphisme de
topos \struck{foncteurs} \ill{} de \ill{} caractérisations universelles de
\ill{})\note{une parenthèse interlinéaire, serrée, dont seuls ces mots se
lisent}
\[
\mathcal{Q}(u) : \mathcal{Q}(\alpha, \beta) \longrightarrow \mathcal{Q}(\lambda, \mu)
\]
avec transitivités, \ill{} rendant commutatif (à isom. can. près) le
diagramme

\begin{tikzcd}
X \arrow[r, "\omega_X"] \arrow[d, "u_0"'] & \mathcal{Q}(\alpha, \beta) \arrow[d, "\mathcal{Q}(u) = v"] \\
Y \arrow[r, "\omega_Y"] & \mathcal{Q}(\lambda, \mu)
\end{tikzcd}

d'ailleurs les foncteurs \ill{} essentiellement \ill{} comme \ill{}
interprète \ill{} le cas \ill{} $\struck{\alpha'}\,X' = X$, $Y' = Y$,
$\struck{\alpha}\,\beta = \mathrm{id}_X$, $\mu = \mathrm{id}_Y$,
$\alpha = F_X$, $\lambda = F_Y$, et on \ill{} commutatif comme \ill{}
dessus …

Revenons au cas général. Si $\mathbb{E} = (E, F_E)$ est un objet de
$\mathcal{Q}(\alpha, \beta)$, i.e. des « \ill{} \ill{} des chgt. de base »,
$\mathcal{Q}(u_0)_*$ et $\mathcal{Q}(u_0)^{*}$ s'explicitent comme au n° 1,
\ill{} \struck{\ill{}} \ill{}\note{les lignes interlinéaires qui suivent, entre
« comme au n° 1 » et « Revenons », sont une variante serrée que le crayon ne
laisse pas lire}

Revenons au cas \ill{} on \ill{} de $\mathcal{Q}(\alpha, \beta)$, i.e. \ill{}
\[
\omega_X^{*}(R^{i}v_*(\mathbb{E})) \longrightarrow R^{i}u_{0*}(\omega_X^{*}(\mathbb{E})) = R^{i}u_{0*}(E) ,
\]
et on aimerait savoir que ce sont des \emph{isoms}, sachant que c'est le
cas pour $i = 0$. On est obligé pratiquement d'établir que \ill{}
objets\note{lu « objets » ; le mot est serré} \emph{injectifs} de
$\mathcal{Q}(\alpha, \beta)$ \ill{} que \ill{} si $\mathbb{E} = (E, F_E)$ est
un faisceau injectif de $\mathcal{Q}(\alpha, \beta)$, alors $E$ est un
faisceau acyclique \add{pour} $u_{0*}$ (SGA 4 V \ill{}). Il suffit pour ceci
que $\alpha$ et $\beta$ commutent aux produits infinis (i.e. $\alpha_!$,
$\beta_!$ existent)

\page{19}

auquel cas il en est de même de $\omega_X$ i.e. $\omega_{X!}$ \ill{}\note{le
feuillet ne porte que cette ligne, qui clôt la page 18 ; la fin de la ligne
est perdue à la marge droite}

\note{la page 20 est une chemise de sa main qui nomme le dossier suivant,
« Fil. de De Rham (à puissances divisées) » ; le dossier lui-même commence au
lot 2}

\end{document}
